\documentclass[a4paper,11pt]{article}
\usepackage[utf8]{inputenc}
\usepackage[T1]{fontenc}
\usepackage[english]{babel}
\usepackage[left=2.5cm,right=2.5cm,top=2cm,bottom=2cm]{geometry}
\usepackage{hyperref}
\usepackage{xcolor}
\usepackage{microtype}
\usepackage{lmodern}

\definecolor{primary}{RGB}{0, 45, 114}
\hypersetup{colorlinks=true, urlcolor=primary}

\begin{document}
\noindent
\textbf{Loïc Aron MBASSI EWOLO} \\
ENSEA -- Double Diplôme Ingénieur \\
\href{mailto:loicaron.mbassiewolo@ensea.fr}{loicaron.mbassiewolo@ensea.fr} \quad (+33) 7 59 52 33 98 \\
\href{https://github.com/Nameless0l}{github.com/Nameless0l} \quad \href{https://mbassiloic.tech}{mbassiloic.tech}

\vspace{1em}
\noindent Cergy, le 22 mars 2026

\vspace{1em}
\noindent
\textbf{Objet : Candidature -- Agents IA Amélioration Continue} \\
\textbf{Référence : Offre d'alternance -- Thales Bordeaux}

\vspace{1.5em}
\noindent Madame, Monsieur,

\vspace{0.8em}
L'amélioration continue exige des systèmes IA capables d'analyser la performance, d'identifier les opportunités et de piloter l'optimisation en boucle fermée. Vous recherchez un ingénieur capable de concevoir des agents IA autonomes pour cette transformation. Mon expérience en orchestration multi-agents (Google ADK avec RAG), en pipeline NLP de recherche (INRIA Saclay) et en leadership technique (Head of AI Cell à ENSPY) me positionne pour co-construire vos agents d'amélioration continue et former les équipes associées.

\vspace{0.8em}
J'achève actuellement mon stage de recherche à INRIA Saclay (équipe TRIBE, avril-août 2026), où je conçois un système d'analyse et de scoring automatisé évaluant la transparence et la conformité de centaines d'applications. Ce travail m'immerge dans la création de pipelines NLP robustes combinant analyse statique, extraction de features textuelles et fusion de signaux pour produire un score de confiance interpétable. Cette expérience m'a montré comment encoder les meilleures pratiques (RGPD, sécurité, transparence) en règles d'évaluation mesurables. Mon système multi-agents agricole (Google ADK, Gemini, RAG) m'a formé à l'orchestration de cinq agents collaboratifs produisant des recommandations multi-critères, un pattern applicable à l'amélioration continue d'une chaîne de production ou d'un service.

\vspace{0.8em}
Le projet Snappy (plateforme temps réel sécurisée, 4ème année ENSPY) m'a exposé à la conception d'architectures DDD (Domain-Driven Design) où les invariants métier sont respectés à travers une machine à états explicite et des transactions atomiques. Cette rigueur architecurale est indispensable pour que les agents IA restent prévisibles et contrôlables en production. En tant que Head of AI Cell à ENSPY (leadership de plus de 50 membres), j'ai coordonné une équipe pluridisciplinaire, piloté des projets concurrents et formé mes pairs aux meilleures pratiques IA. Cette expérience m'a appris comment échelle les initiatives et créer une dynamique d'apprentissage continu.

\vspace{0.8em}
Je suis disponible dès septembre 2026 pour une alternance de 12 mois. Convaincu que l'amélioration continue repose sur des agents IA autonomes, transparents et contrôlables, je suis prêt à contribuer à la transformation de Thales Bordeaux.

\vspace{1.5em}
\noindent Veuillez agréer, Madame, Monsieur, l'expression de mes salutations distinguées.

\vspace{2em}
\noindent \textbf{Loïc Aron MBASSI EWOLO}

\end{document}
